\chapter{Additional files}\label{app:additional}\label{app:noiseinject}\label{app:results}
% ---- TAKEAWAYS ----
% - Every supporting table and figure too long for a main chapter is here, as one numbered Additional file, grouped by the chapter it supports.
% - The Chapter 3 files are the NoiseInject paper's Additional files; ten of the eleven come from the earlier version of the study and must be replaced.
% - The Chapter 5 files let a reader rebuild any hybrid: every model's settings, the models left out per dataset, the size of the search, and what a blueprint saves.
% - The design-experiment file shows where each search default came from. Those numbers are preliminary design runs, not results of the thesis.
% - The Chapter 4 and 6 files hold the full per-dataset tables that those chapters summarise. Most need the full runs.
% ---- END TAKEAWAYS ----
\newcounter{addfile}
\providecommand{\theHaddfile}{addfile.\arabic{addfile}}
\newcommand{\addfile}[1]{\refstepcounter{addfile}\par\medskip\noindent\textbf{Additional file~\theaddfile: #1}\par\nopagebreak\smallskip}

% TODO: review, needs manual review
This appendix holds the tables and figures that support the main chapters but are too long to include in them. They are grouped by the chapter they support and numbered in order as Additional files. The files for Chapter~\ref{chap:robustness} follow the Additional files of the paper that chapter is adapted from.

\section{Noise robustness across molecular representations and model architectures}\label{app:add:robustness}
% ---- NOTES: Noise robustness across molecular representations and model architectures ----
% Takeaway: the supporting tables and figures for Chapter 3, in the NoiseInject paper's grouping.
% - Out of date: every file but the first. They use the noise degradation slope, six noise strategies and Mol2Vec. Chapter 3 now uses the normalised area under the retention curve.
% Result needed: replace the ten out-of-date tables and figures once the Chapter 3 results are updated.
% - Exclusion threshold: clean $R^2 < 0.3$, per replicate on QM9 and per fold on the assay datasets.
% - Change: the excluded-configurations caption and the Chapter 3 Results paragraph (66 configurations) still use the earlier threshold of 0.6.
% Result needed: excluded configurations at clean $R^2 < 0.3$, one row per model, representation and noise condition.
% - The redundancy tables also list five representations Chapter 3 does not use (Morgan, randomised SMILES, SMILES, binary PDV, Mol2Vec).
% - The redundancy and ICC files belong to the correlation filter before the ANOVA; the paper replaced it with model families. Which one to keep is still open.
% - One new file is planned, cited by Chapter 3: base models against their probabilistic counterparts on LogD, Caco-2 efflux and hERG $K_i$, laid out like the QM9 table in Chapter 3.
% - Each entry: median change in AUC_norm, then clean $R^2$ in brackets, counterpart minus base, over 5 folds. One column says whether the pair shares one setting.
% Result needed: that table, on the three assay datasets.
% ---- END NOTES ----

\addfile{Hyperparameters for every model configuration, and the descriptors that make up PDV}
\label{app:noise:sec:hyperparameters_default}\label{app:noise:sec:hyperparameters_tuned}\label{app:noise:sec:pdv_descriptors}
% TODO: review, needs manual review
Table~A lists the default settings of every model configuration, Table~B the tuned settings that replace a default on each dataset, and Table~C the 200 descriptors that make up PDV.
% Thesis appendix item: hyperparameters_default
% Seeded on 2026-09-27 from qsar_qm_models/supplementary/items/hyperparameters_default.tex by thesis_appendix/seed_from_paper.py.
% From here on this copy is edited independently of the paper. Rules: thesis_appendix/README.md.
% The lines below, down to the first blank line, are the PAPER copy's header, kept for
% provenance. They describe the paper file, not this one: in particular a GENERATED
% notice there does not apply here, and this copy may be edited by hand.
% Supplementary item: hyperparameters_default
% GENERATED by scripts/generate_supp_table1.py --write. Do not edit;
% scripts/test_supp_table1.py fails on a hand edit. Rules: supplementary/README.md.

{\small
% TODO: REVIEW MODIFICATION (30 Sep) -- was p{5.0cm}p{5.0cm}p{4.6cm}, 12pt wider than the text; the first column is wider so
% the configuration names fit on one line.
\begin{longtable}{@{}p{5.9cm}p{4.6cm}p{3.6cm}@{}}
\caption{\textbf{Table A: default hyperparameters for every model configuration in the study.}
Every value is read from \texttt{models/model\_defaults.py} at spec version 1.7.0, spec hash \texttt{9249926c38dc}, which is the file both the QM9 pipeline and the assay pipeline load their settings from. Every QM9 results row carries the spec version and the spec hash as two of its columns, so a QM9 table can be traced back to the settings that produced it; the three assay datasets are written by a second pipeline whose results files do not carry those two columns. The configuration list is the model roster the job generator submits. One setting is used for every representation and for every dataset, except where Table B gives a tuned replacement. The support vector machine uses an RBF kernel on every representation, so no result in this study carries a kernel that changes with the representation. Values not held in the shared spec name their source file in the third column.}
\label{af1:hyperparameters}\\
\toprule
\textbf{Configuration} & \textbf{Hyperparameter} & \textbf{Default value} \\
\midrule
\endfirsthead
\toprule
\textbf{Configuration} & \textbf{Hyperparameter} & \textbf{Default value} \\
\midrule
\endhead
\multicolumn{3}{@{}l}{\textit{Shared by every neural network below}} \\
\midrule
Training, all networks & optimizer & adam \\
 & lr & 0.001 \\
 & weight\_decay & 0.0 \\
 & batch\_size & 32 \\
 & epochs & 100 \newline {\footnotesize\raggedright cap; early stopping usually reaches it first} \\
 & patience & 10 \newline {\footnotesize\raggedright epochs without a strictly better mean validation loss} \\
 & restore\_best\_weights & True \newline {\footnotesize\raggedright the reported fit is the best epoch, not the last} \\
 & val\_loss\_reduction & mean \\
 & improvement\_tolerance & 0.0 \\
 & mc\_passes & 100 \newline {\footnotesize\raggedright stochastic forward passes at inference} \\
 & standardise\_targets & True \newline {\footnotesize\raggedright fitted on training labels alone} \\
\midrule
\multicolumn{3}{@{}l}{\textit{Tree and boosting models}} \\
\midrule
RF \newline \footnotesize\texttt{rf} & n\_estimators & 100 \\
 & max\_depth & None \newline {\footnotesize\raggedright no limit} \\
 & min\_samples\_leaf & 5 \\
 & min\_samples\_split & 2 \\
 & max\_features & 0.3 \\
 & bootstrap & True \\
\addlinespace
RF300 \newline \footnotesize\texttt{rf300} & n\_estimators & 300 \\
 & max\_depth & None \newline {\footnotesize\raggedright no limit} \\
 & min\_samples\_leaf & 5 \\
 & min\_samples\_split & 2 \\
 & max\_features & 0.3 \\
 & bootstrap & True \\
\addlinespace
XGBoost \newline \footnotesize\texttt{xgboost} & n\_estimators & 100 \\
 & max\_depth & 6 \\
 & learning\_rate & 0.1 \\
 & subsample & 1.0 \\
 & colsample\_bytree & 1.0 \\
 & colsample\_bylevel & 1.0 \\
 & min\_child\_weight & 1 \\
 & gamma & 0.0 \\
 & reg\_alpha & 0.0 \\
 & reg\_lambda & 1.0 \\
\addlinespace
LightGBM \newline \footnotesize\texttt{lgb} & n\_estimators & 100 \\
 & learning\_rate & 0.1 \\
 & num\_leaves & 31 \\
 & max\_depth & $-1$ \newline {\footnotesize\raggedright no limit} \\
 & subsample & 1.0 \\
 & colsample\_bytree & 1.0 \\
 & min\_child\_samples & 20 \\
 & reg\_alpha & 0.0 \\
 & reg\_lambda & 0.0 \\
\addlinespace
NGBoost \newline \footnotesize\texttt{ngboost} & n\_estimators & 500 \\
 & learning\_rate & 0.01 \\
 & natural\_gradient & True \\
 & dist & Normal \newline {\footnotesize\raggedright ngboost.distns.Normal} \\
 & score & MLE \newline {\footnotesize\raggedright ngboost.scores.MLE, the log scoring rule} \\
 & minibatch\_frac & 1.0 \\
 & col\_sample & 1.0 \\
 & base\_max\_depth & 3 \newline {\footnotesize\raggedright depth of one boosting stage} \\
 & base\_criterion & friedman\_mse \\
 & early\_stopping\_rounds & 50 \newline {\footnotesize\raggedright stages without a better validation log likelihood} \\
 & use\_best\_iteration & True \newline {\footnotesize\raggedright predictions use the best stage, not every stage fitted} \\
 & ngboost\_ensemble\_seeds & 1 \newline {\footnotesize\raggedright 1 = one fit, no seed ensemble} \\
\addlinespace
QRF \newline \footnotesize\texttt{qrf} & n\_estimators & 300 \\
 & max\_depth & None \newline {\footnotesize\raggedright no limit} \\
 & min\_samples\_leaf & 5 \\
 & min\_samples\_split & 2 \\
 & max\_features & 0.3 \\
 & bootstrap & True \\
\midrule
\multicolumn{3}{@{}l}{\textit{Kernel models}} \\
\midrule
SVM \newline \footnotesize\texttt{svm} & C & 1.0 \\
 & gamma & scale \\
 & kernel & rbf \\
\addlinespace
GP \newline \footnotesize\texttt{gauche\_rbf} & kernel & rbf \newline {\footnotesize\raggedright set by this job, not by the shared spec} \\
 & likelihood\_noise & 0.001 \newline {\footnotesize\raggedright starting value; it is learned} \\
 & outputscale & 1.0 \\
 & apply\_outputscale & True \\
 & max\_train\_n & 5000 \newline {\footnotesize\raggedright molecules the process is fitted on; an exact Gaussian process is cubic in this. The test set is untouched and the number actually fitted is written into every results row} \\
 & fallback\_adam\_lr & 0.1 \newline {\footnotesize\raggedright used only when the standard fitter refuses the model} \\
 & fallback\_adam\_iters & 100 \newline {\footnotesize\raggedright used only when the standard fitter refuses the model} \\
 & single\_thread\_fit & False \newline {\footnotesize\raggedright a threading workaround, off since the environment check found one runtime} \\
 & init\_lengthscale\_from\_data & True \newline {\footnotesize\raggedright the RBF length scale starts at the median distance between training molecules, not at the library default} \\
 & lengthscale\_probe\_n & 500 \newline {\footnotesize\raggedright molecules sampled to estimate that median} \\
 & collapse\_fraction & 0.05 \newline {\footnotesize\raggedright a fit whose predictions vary by less than this fraction of the training label spread is recorded as collapsed rather than scored} \\
\addlinespace
GP-Tanimoto \newline \footnotesize\texttt{gauche} \newline \footnotesize run on ecfp4 only & kernel & tanimoto \newline {\footnotesize\raggedright set by this job, not by the shared spec} \\
 & likelihood\_noise & 0.001 \newline {\footnotesize\raggedright starting value; it is learned} \\
 & outputscale & 1.0 \\
 & apply\_outputscale & True \\
 & max\_train\_n & 5000 \newline {\footnotesize\raggedright molecules the process is fitted on; an exact Gaussian process is cubic in this. The test set is untouched and the number actually fitted is written into every results row} \\
 & fallback\_adam\_lr & 0.1 \newline {\footnotesize\raggedright used only when the standard fitter refuses the model} \\
 & fallback\_adam\_iters & 100 \newline {\footnotesize\raggedright used only when the standard fitter refuses the model} \\
 & single\_thread\_fit & False \newline {\footnotesize\raggedright a threading workaround, off since the environment check found one runtime} \\
 & init\_lengthscale\_from\_data & True \newline {\footnotesize\raggedright the RBF length scale starts at the median distance between training molecules, not at the library default} \\
 & lengthscale\_probe\_n & 500 \newline {\footnotesize\raggedright molecules sampled to estimate that median} \\
 & collapse\_fraction & 0.05 \newline {\footnotesize\raggedright a fit whose predictions vary by less than this fraction of the training label spread is recorded as collapsed rather than scored} \\
\addlinespace
GP-Hetero \newline \footnotesize\texttt{heteroscedastic\_gp} & kernel & rbf \newline {\footnotesize\raggedright set by this job, not by the shared spec} \\
 & likelihood\_noise & 0.001 \newline {\footnotesize\raggedright starting value; it is learned} \\
 & outputscale & 1.0 \\
 & apply\_outputscale & True \\
 & max\_train\_n & 5000 \newline {\footnotesize\raggedright molecules the process is fitted on; an exact Gaussian process is cubic in this. The test set is untouched and the number actually fitted is written into every results row} \\
 & fallback\_adam\_lr & 0.1 \newline {\footnotesize\raggedright used only when the standard fitter refuses the model} \\
 & fallback\_adam\_iters & 100 \newline {\footnotesize\raggedright used only when the standard fitter refuses the model} \\
 & single\_thread\_fit & False \newline {\footnotesize\raggedright a threading workaround, off since the environment check found one runtime} \\
 & init\_lengthscale\_from\_data & True \newline {\footnotesize\raggedright the RBF length scale starts at the median distance between training molecules, not at the library default} \\
 & lengthscale\_probe\_n & 500 \newline {\footnotesize\raggedright molecules sampled to estimate that median} \\
 & collapse\_fraction & 0.05 \newline {\footnotesize\raggedright a fit whose predictions vary by less than this fraction of the training label spread is recorded as collapsed rather than scored} \\
 & noise network hidden sizes & [64, 64] \newline {\footnotesize\raggedright models/models.py, not in the shared spec} \\
 & noise network output & Softplus, plus $10^{-4}$ \newline {\footnotesize\raggedright models/models.py, not in the shared spec} \\
 & Adam learning rate, process & 0.1 \newline {\footnotesize\raggedright models/models.py, not in the shared spec} \\
 & Adam learning rate, noise network & 0.001 \newline {\footnotesize\raggedright models/models.py, not in the shared spec} \\
 & noise network epochs & 100 \newline {\footnotesize\raggedright the shared epoch cap} \\
\midrule
\multicolumn{3}{@{}l}{\textit{Neural networks}} \\
\midrule
NN-$\alpha$ \newline \footnotesize\texttt{dnn} & hidden\_sizes & [128, 64] \\
 & dropout\_rate & 0.2 \\
 & activation & relu \\
\addlinespace
NN-$\beta$ \newline \footnotesize\texttt{mlp} & hidden\_size & 128 \\
 & num\_hidden\_layers & 2 \\
 & dropout\_rate & 0.2 \\
 & activation & relu \\
\addlinespace
BNN-$\alpha$ \newline \footnotesize\texttt{dnn\_bnn\_full} & hidden\_sizes & [128, 64] \\
 & dropout\_rate & 0.2 \\
 & activation & relu \\
 & bnn\_prior\_mu & 0.0 \\
 & bnn\_prior\_sigma & 0.1 \\
 & bnn\_kl\_weight & elbo \newline {\footnotesize\raggedright 1 / number of training molecules} \\
\addlinespace
BNN-$\beta$ \newline \footnotesize\texttt{mlp\_bnn\_full} & hidden\_size & 128 \\
 & num\_hidden\_layers & 2 \\
 & dropout\_rate & 0.2 \\
 & activation & relu \\
 & bnn\_prior\_mu & 0.0 \\
 & bnn\_prior\_sigma & 0.1 \\
 & bnn\_kl\_weight & elbo \newline {\footnotesize\raggedright 1 / number of training molecules} \\
\addlinespace
VBLL-$\alpha$ \newline \footnotesize\texttt{dnn\_bnn\_full\_variational} & hidden\_sizes & [128, 64] \\
 & dropout\_rate & 0.2 \\
 & activation & relu \\
 & vbll\_prior\_mu & 0.0 \\
 & vbll\_prior\_sigma & 1.0 \\
 & vbll\_init\_log\_sigma & $-3.0$ \\
 & vbll\_init\_log\_noise\_var & $-2.0$ \\
\addlinespace
VBLL-$\beta$ \newline \footnotesize\texttt{mlp\_bnn\_full\_variational} & hidden\_size & 128 \\
 & num\_hidden\_layers & 2 \\
 & dropout\_rate & 0.2 \\
 & activation & relu \\
 & vbll\_prior\_mu & 0.0 \\
 & vbll\_prior\_sigma & 1.0 \\
 & vbll\_init\_log\_sigma & $-3.0$ \\
 & vbll\_init\_log\_noise\_var & $-2.0$ \\
\addlinespace
VBLL-$\alpha$ (noise head) \newline \footnotesize\texttt{dnn\_bnn\_full\_variational\_hetero} & hidden\_sizes & [128, 64] \\
 & dropout\_rate & 0.2 \\
 & activation & relu \\
 & vbll\_prior\_mu & 0.0 \\
 & vbll\_prior\_sigma & 1.0 \\
 & vbll\_init\_log\_sigma & $-3.0$ \\
 & vbll\_init\_log\_noise\_var & $-2.0$ \\
 & observation noise & predicted per molecule \newline {\footnotesize\raggedright a variational noise head on the same features, started at the single-number value above} \\
\addlinespace
VBLL-$\beta$ (noise head) \newline \footnotesize\texttt{mlp\_bnn\_full\_variational\_hetero} & hidden\_size & 128 \\
 & num\_hidden\_layers & 2 \\
 & dropout\_rate & 0.2 \\
 & activation & relu \\
 & vbll\_prior\_mu & 0.0 \\
 & vbll\_prior\_sigma & 1.0 \\
 & vbll\_init\_log\_sigma & $-3.0$ \\
 & vbll\_init\_log\_noise\_var & $-2.0$ \\
 & observation noise & predicted per molecule \newline {\footnotesize\raggedright a variational noise head on the same features, started at the single-number value above} \\
\addlinespace
BNN-$\alpha$ (variance head) \newline \footnotesize\texttt{dnn\_bnn\_full\_mve} & hidden\_sizes & [128, 64] \\
 & dropout\_rate & 0.2 \\
 & activation & relu \\
 & bnn\_prior\_mu & 0.0 \\
 & bnn\_prior\_sigma & 0.1 \\
 & bnn\_kl\_weight & elbo \newline {\footnotesize\raggedright 1 / number of training molecules} \\
 & output head & value and log variance \newline {\footnotesize\raggedright two outputs} \\
 & loss & Gaussian negative log likelihood \\
 & log variance clamp & $[-10, 10]$ \newline {\footnotesize\raggedright scripts/loss\_functions.py, not in the shared spec} \\
\addlinespace
BNN-$\beta$ (variance head) \newline \footnotesize\texttt{mlp\_bnn\_full\_mve} & hidden\_size & 128 \\
 & num\_hidden\_layers & 2 \\
 & dropout\_rate & 0.2 \\
 & activation & relu \\
 & bnn\_prior\_mu & 0.0 \\
 & bnn\_prior\_sigma & 0.1 \\
 & bnn\_kl\_weight & elbo \newline {\footnotesize\raggedright 1 / number of training molecules} \\
 & output head & value and log variance \newline {\footnotesize\raggedright two outputs} \\
 & loss & Gaussian negative log likelihood \\
 & log variance clamp & $[-10, 10]$ \newline {\footnotesize\raggedright scripts/loss\_functions.py, not in the shared spec} \\
\bottomrule
\end{longtable}
}

% Thesis appendix item: hyperparameters_tuned
% Seeded on 2026-09-27 from qsar_qm_models/supplementary/items/hyperparameters_tuned.tex by thesis_appendix/seed_from_paper.py.
% From here on this copy is edited independently of the paper. Rules: thesis_appendix/README.md.
% The lines below, down to the first blank line, are the PAPER copy's header, kept for
% provenance. They describe the paper file, not this one: in particular a GENERATED
% notice there does not apply here, and this copy may be edited by hand.
% Supplementary item: hyperparameters_tuned
% GENERATED by scripts/generate_supp_table1.py --write. Do not edit;
% scripts/test_supp_table1.py fails on a hand edit. Rules: supplementary/README.md.

{\small
\begin{longtable}{@{}llll@{}}
\caption{\textbf{The tuned settings that replace a default, and the four datasets they apply to.}
A setting is chosen once per model per dataset, not per representation, and is then used for all six representations: a two-way comparison across model and representation cannot be read if each cell carries its own hyperparameters. Only the Bayesian and variational networks are tuned; every other configuration in Table~\ref{af1:hyperparameters} trains at the default. A model with no row for a dataset kept the default on that dataset. Learning rates are printed to three significant figures; the source files hold the full value. Source files: \texttt{results/master\_tuned\_hyperparameters.json} for QM9 and \texttt{results/master\_tuned\_hyperparameters\_lab.json} for the three assay datasets.}
\label{af1:tuned}\\
\toprule
\textbf{Dataset} & \textbf{Configuration} & \textbf{Hyperparameter} & \textbf{Tuned value} \\
\midrule
\endfirsthead
\toprule
\textbf{Dataset} & \textbf{Configuration} & \textbf{Hyperparameter} & \textbf{Tuned value} \\
\midrule
\endhead
QM9 & NN-$\alpha$ & activation & tanh \\
 &  & hidden\_size1 & 64 \\
 &  & hidden\_size2 & 64 \\
\addlinespace
QM9 & BNN-$\alpha$ & activation & tanh \\
 &  & hidden\_size1 & 64 \\
 &  & hidden\_size2 & 64 \\
\addlinespace
QM9 & VBLL-$\alpha$ & activation & tanh \\
 &  & hidden\_size1 & 64 \\
 &  & hidden\_size2 & 64 \\
\addlinespace
QM9 & NN-$\beta$ & dropout\_rate & 0.401 \\
 &  & hidden\_size & 64 \\
 &  & lr & 0.00587 \\
 &  & num\_hidden\_layers & 3 \\
\addlinespace
QM9 & BNN-$\beta$ & dropout\_rate & 0.401 \\
 &  & hidden\_size & 64 \\
 &  & lr & 0.00587 \\
 &  & num\_hidden\_layers & 3 \\
\addlinespace
QM9 & VBLL-$\beta$ & dropout\_rate & 0.357 \\
 &  & hidden\_size & 64 \\
 &  & lr & 0.00119 \\
 &  & num\_hidden\_layers & 1 \\
\addlinespace
Caco-2 & NN-$\alpha$ & activation & tanh \\
 &  & hidden\_size1 & 64 \\
 &  & hidden\_size2 & 64 \\
\addlinespace
Caco-2 & BNN-$\alpha$ & activation & tanh \\
 &  & hidden\_size1 & 64 \\
 &  & hidden\_size2 & 64 \\
\addlinespace
Caco-2 & VBLL-$\alpha$ & activation & relu \\
 &  & hidden\_size1 & 64 \\
 &  & hidden\_size2 & 32 \\
\addlinespace
Caco-2 & NN-$\beta$ & dropout\_rate & 0.401 \\
 &  & hidden\_size & 64 \\
 &  & lr & 0.00587 \\
 &  & num\_hidden\_layers & 3 \\
\addlinespace
Caco-2 & BNN-$\beta$ & dropout\_rate & 0.401 \\
 &  & hidden\_size & 64 \\
 &  & lr & 0.00587 \\
 &  & num\_hidden\_layers & 3 \\
\addlinespace
hERG K$_i$ & NN-$\alpha$ & activation & tanh \\
 &  & hidden\_size1 & 64 \\
 &  & hidden\_size2 & 64 \\
\addlinespace
hERG K$_i$ & BNN-$\alpha$ & activation & tanh \\
 &  & hidden\_size1 & 64 \\
 &  & hidden\_size2 & 64 \\
\addlinespace
hERG K$_i$ & VBLL-$\alpha$ & activation & tanh \\
 &  & hidden\_size1 & 64 \\
 &  & hidden\_size2 & 64 \\
\addlinespace
hERG K$_i$ & NN-$\beta$ & dropout\_rate & 0.401 \\
 &  & hidden\_size & 64 \\
 &  & lr & 0.00587 \\
 &  & num\_hidden\_layers & 3 \\
\addlinespace
hERG K$_i$ & BNN-$\beta$ & dropout\_rate & 0.401 \\
 &  & hidden\_size & 64 \\
 &  & lr & 0.00587 \\
 &  & num\_hidden\_layers & 3 \\
\addlinespace
LogD & NN-$\alpha$ & activation & tanh \\
 &  & hidden\_size1 & 64 \\
 &  & hidden\_size2 & 64 \\
\addlinespace
LogD & BNN-$\alpha$ & activation & tanh \\
 &  & hidden\_size1 & 64 \\
 &  & hidden\_size2 & 64 \\
\addlinespace
LogD & VBLL-$\alpha$ & activation & tanh \\
 &  & hidden\_size1 & 256 \\
 &  & hidden\_size2 & 32 \\
\addlinespace
LogD & NN-$\beta$ & dropout\_rate & 0.401 \\
 &  & hidden\_size & 64 \\
 &  & lr & 0.00587 \\
 &  & num\_hidden\_layers & 3 \\
\addlinespace
LogD & BNN-$\beta$ & dropout\_rate & 0.401 \\
 &  & hidden\_size & 64 \\
 &  & lr & 0.00587 \\
 &  & num\_hidden\_layers & 3 \\
\addlinespace
LogD & VBLL-$\beta$ & dropout\_rate & 0.079 \\
 &  & hidden\_size & 128 \\
 &  & lr & 0.00011 \\
 &  & num\_hidden\_layers & 2 \\
\bottomrule
\end{longtable}
}

% Thesis appendix item: pdv_descriptors
% Seeded on 2026-09-27 from qsar_qm_models/supplementary/items/pdv_descriptors.tex by thesis_appendix/seed_from_paper.py.
% From here on this copy is edited independently of the paper. Rules: thesis_appendix/README.md.
% The lines below, down to the first blank line, are the PAPER copy's header, kept for
% provenance. They describe the paper file, not this one: in particular a GENERATED
% notice there does not apply here, and this copy may be edited by hand.
% Supplementary item: pdv_descriptors
% GENERATED by scripts/generate_supp_table1.py --write. Do not edit;
% scripts/test_supp_table1.py fails on a hand edit. Rules: supplementary/README.md.

% [2026-10-10 fix] kept \scriptsize (30 Sep change; \small was 70pt wider than the text); additional_files_body.tex also sets \small to \scriptsize around this \input
{\scriptsize
\begin{longtable}{@{}llll@{}}
\caption{\textbf{The 200 descriptors that make up the PDV representation.}
One entry is one feature of the vector, in the order the pipeline computes them. The list is fixed in the code rather than taken from RDKit's current default set, which grows between releases, so the same 200 features are built on every dataset and in every run. They were computed with RDKit 2022.09.5. Descriptors that return a non-finite value on a molecule have that value replaced by zero before the standardization constants are computed. Source: \texttt{DEFAULT\_DESCRIPTOR\_LIST} in \texttt{scripts/process\_and\_train.py}.}
\label{af1:pdv}\\
\toprule
\endfirsthead
\toprule
\endhead
BalabanJ & NOCount & SlogP\_VSA7 & fr\_bicyclic \\
BertzCT & NumAliphaticCarbocycles & SlogP\_VSA8 & fr\_diazo \\
Chi0 & NumAliphaticHeterocycles & SlogP\_VSA9 & fr\_dihydropyridine \\
Chi0n & NumAliphaticRings & TPSA & fr\_epoxide \\
Chi0v & NumAromaticCarbocycles & VSA\_EState1 & fr\_ester \\
Chi1 & NumAromaticHeterocycles & VSA\_EState10 & fr\_ether \\
Chi1n & NumAromaticRings & VSA\_EState2 & fr\_furan \\
Chi1v & NumHAcceptors & VSA\_EState3 & fr\_guanido \\
Chi2n & NumHDonors & VSA\_EState4 & fr\_halogen \\
Chi2v & NumHeteroatoms & VSA\_EState5 & fr\_hdrzine \\
Chi3n & NumRadicalElectrons & VSA\_EState6 & fr\_hdrzone \\
Chi3v & NumRotatableBonds & VSA\_EState7 & fr\_imidazole \\
Chi4n & NumSaturatedCarbocycles & VSA\_EState8 & fr\_imide \\
Chi4v & NumSaturatedHeterocycles & VSA\_EState9 & fr\_isocyan \\
EState\_VSA1 & NumSaturatedRings & fr\_Al\_COO & fr\_isothiocyan \\
EState\_VSA10 & NumValenceElectrons & fr\_Al\_OH & fr\_ketone \\
EState\_VSA11 & PEOE\_VSA1 & fr\_Al\_OH\_noTert & fr\_ketone\_Topliss \\
EState\_VSA2 & PEOE\_VSA10 & fr\_ArN & fr\_lactam \\
EState\_VSA3 & PEOE\_VSA11 & fr\_Ar\_COO & fr\_lactone \\
EState\_VSA4 & PEOE\_VSA12 & fr\_Ar\_N & fr\_methoxy \\
EState\_VSA5 & PEOE\_VSA13 & fr\_Ar\_NH & fr\_morpholine \\
EState\_VSA6 & PEOE\_VSA14 & fr\_Ar\_OH & fr\_nitrile \\
EState\_VSA7 & PEOE\_VSA2 & fr\_COO & fr\_nitro \\
EState\_VSA8 & PEOE\_VSA3 & fr\_COO2 & fr\_nitro\_arom \\
EState\_VSA9 & PEOE\_VSA4 & fr\_C\_O & fr\_nitro\_arom\_nonortho \\
ExactMolWt & PEOE\_VSA5 & fr\_C\_O\_noCOO & fr\_nitroso \\
FpDensityMorgan1 & PEOE\_VSA6 & fr\_C\_S & fr\_oxazole \\
FpDensityMorgan2 & PEOE\_VSA7 & fr\_HOCCN & fr\_oxime \\
FpDensityMorgan3 & PEOE\_VSA8 & fr\_Imine & fr\_para\_hydroxylation \\
FractionCSP3 & PEOE\_VSA9 & fr\_NH0 & fr\_phenol \\
HallKierAlpha & RingCount & fr\_NH1 & fr\_phenol\_noOrthoHbond \\
HeavyAtomCount & SMR\_VSA1 & fr\_NH2 & fr\_phos\_acid \\
HeavyAtomMolWt & SMR\_VSA10 & fr\_N\_O & fr\_phos\_ester \\
Ipc & SMR\_VSA2 & fr\_Ndealkylation1 & fr\_piperdine \\
Kappa1 & SMR\_VSA3 & fr\_Ndealkylation2 & fr\_piperzine \\
Kappa2 & SMR\_VSA4 & fr\_Nhpyrrole & fr\_priamide \\
Kappa3 & SMR\_VSA5 & fr\_SH & fr\_prisulfonamd \\
LabuteASA & SMR\_VSA6 & fr\_aldehyde & fr\_pyridine \\
MaxAbsEStateIndex & SMR\_VSA7 & fr\_alkyl\_carbamate & fr\_quatN \\
MaxAbsPartialCharge & SMR\_VSA8 & fr\_alkyl\_halide & fr\_sulfide \\
MaxEStateIndex & SMR\_VSA9 & fr\_allylic\_oxid & fr\_sulfonamd \\
MaxPartialCharge & SlogP\_VSA1 & fr\_amide & fr\_sulfone \\
MinAbsEStateIndex & SlogP\_VSA10 & fr\_amidine & fr\_term\_acetylene \\
MinAbsPartialCharge & SlogP\_VSA11 & fr\_aniline & fr\_tetrazole \\
MinEStateIndex & SlogP\_VSA12 & fr\_aryl\_methyl & fr\_thiazole \\
MinPartialCharge & SlogP\_VSA2 & fr\_azide & fr\_thiocyan \\
MolLogP & SlogP\_VSA3 & fr\_azo & fr\_thiophene \\
MolMR & SlogP\_VSA4 & fr\_barbitur & fr\_unbrch\_alkane \\
MolWt & SlogP\_VSA5 & fr\_benzene & fr\_urea \\
NHOHCount & SlogP\_VSA6 & fr\_benzodiazepine & qed \\
\bottomrule
\end{longtable}
}


\addfile{Configurations excluded from the robustness analysis}
\label{app:noise:sec:excluded_configurations}
% TODO: review, needs manual review
For each model and representation on the QM9 HOMO--LUMO gap, the number of noise strategies for which the configuration was excluded because its clean-label $R^2$ was too low.
% Thesis appendix item: excluded_configurations
% Seeded on 2026-09-27 from qsar_qm_models/supplementary/items/excluded_configurations.tex by thesis_appendix/seed_from_paper.py.
% From here on this copy is edited independently of the paper. Rules: thesis_appendix/README.md.
% The lines below, down to the first blank line, are the PAPER copy's header, kept for
% provenance. They describe the paper file, not this one: in particular a GENERATED
% notice there does not apply here, and this copy may be edited by hand.
% Supplementary item: excluded_configurations
% Split out of additional_files.tex on 2026-09-27, where it was Additional file 5 (Excluded configurations (was Supp Table S2)).
% Content unchanged apart from: the heading line removed (the build writes it), the old
% 'Table S*/Figure S*' caption prefix replaced by \SuppFile{}, and the label renamed to supp:excluded_configurations.
% Rules: supplementary/README.md.

\begin{table}[H]
\centering
\caption{\textbf{Configurations excluded from robustness analysis due to baseline R$^2 \leq 0.6$ on the QM9 HOMO--LUMO gap, summarized by model and representation.}
Each cell shows the number of noise strategies (out of 6) for which the configuration was excluded.}
\label{app:noise:excluded_configurations}
\small
\begin{tabular}{lcccc}
\toprule
\textbf{Model} & \textbf{Mol2Vec} & \textbf{MHGGNN} & \textbf{SMILES} & \textbf{Rand.~SMILES} \\
\midrule
BNN-$\alpha$ & 6 & 6 & --- & --- \\
BNN-$\beta$ & 6 & 6 & --- & --- \\
VBLL-$\alpha$ & 6 & 6 & --- & --- \\
VBLL-$\beta$ & 6 & 6 & --- & --- \\
RF & --- & --- & --- & 4 \\
QRF & --- & --- & --- & 4 \\
NGBoost & --- & --- & 6 & 4 \\
\bottomrule
\end{tabular}
\end{table}


\addfile{Spearman correlations between the robustness profiles of the models}
\label{app:noise:sec:model_redundancy}
% TODO: review, needs manual review
Pairs of models whose noise degradation slope profiles on the QM9 HOMO--LUMO gap correlate at $\rho \geq 0.95$, and which model of each pair was left out of the ANOVA.
% Thesis appendix item: model_redundancy
% Seeded on 2026-09-27 from qsar_qm_models/supplementary/items/model_redundancy.tex by thesis_appendix/seed_from_paper.py.
% From here on this copy is edited independently of the paper. Rules: thesis_appendix/README.md.
% The lines below, down to the first blank line, are the PAPER copy's header, kept for
% provenance. They describe the paper file, not this one: in particular a GENERATED
% notice there does not apply here, and this copy may be edited by hand.
% Supplementary item: model_redundancy
% Split out of additional_files.tex on 2026-09-27, where it was Additional file 2 (Pairwise model NDS correlations (was Supp Table S3)).
% Content unchanged apart from: the heading line removed (the build writes it), the old
% 'Table S*/Figure S*' caption prefix replaced by \SuppFile{}, and the label renamed to supp:model_redundancy.
% Rules: supplementary/README.md.

{\scriptsize
\begin{longtable}{llccc}
\caption{\textbf{Pairwise Spearman rank correlations between model noise degradation slope (NDS) profiles on the QM9 HOMO--LUMO gap ($\rho \geq 0.95$).}
``Model A'' and ``Model B'' label each pair, both treated symmetrically. Checkmarks indicate pairs where one model was excluded from the ANOVA due to near-redundancy. Gauche (RBF) pairs omitted (all $\rho = 1.0$, $n = 6$, PDV-only artifact).}
\label{app:noise:model_redundancy}\\
\toprule
\textbf{Model A} & \textbf{Model B} & \textbf{$n$} & \textbf{$\rho$} & \textbf{Excluded} \\
\midrule
\endfirsthead
\toprule
\textbf{Model A} & \textbf{Model B} & \textbf{$n$} & \textbf{$\rho$} & \textbf{Excluded} \\
\midrule
\endhead
VBLL-$\beta$ & NGBoost & 18 & 0.998 &  \\
LightGBM & XGBoost & 46 & 0.998 &  \\
Gauche & VBLL-$\beta$ & 18 & 0.996 & \checkmark \\
QRF & RF & 40 & 0.995 & \checkmark \\
VBLL-$\alpha$ & LightGBM & 24 & 0.992 &  \\
RF & XGBoost & 46 & 0.989 &  \\
NGBoost & XGBoost & 40 & 0.988 &  \\
NGBoost & RF & 40 & 0.987 &  \\
LightGBM & RF & 46 & 0.986 &  \\
VBLL-$\beta$ & SVM & 24 & 0.986 &  \\
QRF & XGBoost & 40 & 0.986 & \checkmark \\
NGBoost & QRF & 34 & 0.986 & \checkmark \\
VBLL-$\alpha$ & NGBoost & 18 & 0.986 &  \\
VBLL-$\alpha$ & XGBoost & 24 & 0.984 &  \\
LightGBM & NGBoost & 40 & 0.984 &  \\
LightGBM & QRF & 40 & 0.984 & \checkmark \\
VBLL-$\alpha$ & VBLL-$\beta$ & 24 & 0.983 &  \\
NN-$\alpha$ & NN-$\beta$ & 52 & 0.983 &  \\
LightGBM & SVM & 46 & 0.982 &  \\
BNN-$\beta$ & SVM & 34 & 0.981 &  \\
VBLL-$\alpha$ & SVM & 24 & 0.979 &  \\
BNN-$\alpha$ & VBLL-$\beta$ & 24 & 0.979 &  \\
VBLL-$\alpha$ & QRF & 24 & 0.978 & \checkmark \\
SVM & XGBoost & 46 & 0.978 &  \\
BNN-$\alpha$ & BNN-$\beta$ & 36 & 0.977 &  \\
BNN-$\alpha$ & SVM & 34 & 0.976 &  \\
VBLL-$\alpha$ & Gauche & 18 & 0.975 & \checkmark \\
VBLL-$\alpha$ & RF & 24 & 0.975 &  \\
QRF & SVM & 40 & 0.974 & \checkmark \\
BNN-$\alpha$ & VBLL-$\alpha$ & 24 & 0.974 &  \\
LightGBM & VBLL-$\beta$ & 24 & 0.973 &  \\
VBLL-$\alpha$ & BNN-$\beta$ & 24 & 0.972 &  \\
Gauche & SVM & 40 & 0.972 & \checkmark \\
Gauche & XGBoost & 44 & 0.970 & \checkmark \\
NGBoost & SVM & 40 & 0.970 &  \\
NN-$\alpha$ & BNN-$\beta$ & 36 & 0.970 &  \\
Gauche & LightGBM & 42 & 0.969 & \checkmark \\
RF & SVM & 46 & 0.969 &  \\
BNN-$\beta$ & VBLL-$\beta$ & 24 & 0.968 &  \\
VBLL-$\beta$ & QRF & 24 & 0.966 & \checkmark \\
NN-$\alpha$ & BNN-$\alpha$ & 36 & 0.965 &  \\
VBLL-$\beta$ & XGBoost & 24 & 0.965 &  \\
NN-$\beta$ & SVM & 46 & 0.962 &  \\
NN-$\alpha$ & SVM & 46 & 0.960 &  \\
Gauche & NGBoost & 34 & 0.959 & \checkmark \\
LightGBM & BNN-$\beta$ & 36 & 0.959 &  \\
BNN-$\alpha$ & NN-$\beta$ & 36 & 0.958 &  \\
BNN-$\alpha$ & QRF & 28 & 0.956 & \checkmark \\
Gauche & QRF & 34 & 0.956 & \checkmark \\
VBLL-$\beta$ & RF & 24 & 0.955 &  \\
Gauche & RF & 40 & 0.954 & \checkmark \\
BNN-$\beta$ & XGBoost & 34 & 0.953 &  \\
BNN-$\beta$ & QRF & 28 & 0.952 & \checkmark \\
NN-$\alpha$ & VBLL-$\beta$ & 24 & 0.952 &  \\
NN-$\beta$ & BNN-$\beta$ & 36 & 0.952 &  \\
BNN-$\alpha$ & LightGBM & 36 & 0.951 &  \\
NN-$\beta$ & VBLL-$\beta$ & 24 & 0.950 &  \\
BNN-$\alpha$ & Gauche & 30 & 0.950 & \checkmark \\
\bottomrule
\end{longtable}}


\addfile{Spearman correlations between the robustness profiles of the representations}
\label{app:noise:sec:representation_redundancy}
% TODO: review, needs manual review
Pairwise correlations between the noise degradation slope profiles of the representations on the QM9 HOMO--LUMO gap, and which representation of each pair was left out of the ANOVA.
% Thesis appendix item: representation_redundancy
% Seeded on 2026-09-27 from qsar_qm_models/supplementary/items/representation_redundancy.tex by thesis_appendix/seed_from_paper.py.
% From here on this copy is edited independently of the paper. Rules: thesis_appendix/README.md.
% The lines below, down to the first blank line, are the PAPER copy's header, kept for
% provenance. They describe the paper file, not this one: in particular a GENERATED
% notice there does not apply here, and this copy may be edited by hand.
% Supplementary item: representation_redundancy
% Split out of additional_files.tex on 2026-09-27, where it was Additional file 3 (Pairwise representation NDS correlations (was Supp Table S4)).
% Content unchanged apart from: the heading line removed (the build writes it), the old
% 'Table S*/Figure S*' caption prefix replaced by \SuppFile{}, and the label renamed to supp:representation_redundancy.
% Rules: supplementary/README.md.

{\scriptsize
\begin{longtable}{llccc}
\caption{\textbf{Pairwise Spearman rank correlations between representation noise degradation slope (NDS) profiles on the QM9 HOMO--LUMO gap.}
``Rep A'' and ``Rep B'' label each pair of molecular representations, both treated symmetrically. Checkmarks indicate pairs where one representation was excluded from the ANOVA.}
\label{app:noise:representation_redundancy}\\
\toprule
\textbf{Rep.\ A} & \textbf{Rep.\ B} & \textbf{$n$} & \textbf{$\rho$} & \textbf{Excluded} \\
\midrule
\endfirsthead
\toprule
\textbf{Rep.\ A} & \textbf{Rep.\ B} & \textbf{$n$} & \textbf{$\rho$} & \textbf{Excluded} \\
\midrule
\endhead
ECFP4 & Morgan & 60 & 0.995 & \checkmark \\
MHGGNN & Morgan & 48 & 0.994 & \checkmark \\
Rand.~SMILES & SMILES & 16 & 0.994 & \checkmark \\
Rand.~SMILES & SNS & 16 & 0.994 & \checkmark \\
ECFP4 & MHGGNN & 54 & 0.993 &  \\
ECFP4 & SNS & 56 & 0.992 & \checkmark \\
Morgan & SNS & 52 & 0.992 & \checkmark \\
MHGGNN & SNS & 44 & 0.991 & \checkmark \\
PDV & ECFP4 & 72 & 0.990 &  \\
PDV & MHGGNN & 48 & 0.989 &  \\
SMILES & SNS & 52 & 0.988 & \checkmark \\
Morgan & Rand.~SMILES & 16 & 0.988 & \checkmark \\
Morgan & SMILES & 54 & 0.988 & \checkmark \\
PDV & Morgan & 54 & 0.986 & \checkmark \\
PDV & Rand.~SMILES & 12 & 0.986 & \checkmark \\
PDV & SNS & 50 & 0.984 & \checkmark \\
ECFP4 & SMILES & 72 & 0.984 &  \\
Mol2Vec & Morgan & 48 & 0.982 & \checkmark \\
MHGGNN & Mol2Vec & 54 & 0.982 &  \\
ECFP4 & Mol2Vec & 54 & 0.982 &  \\
PDV & PDV (binary) & 72 & 0.980 & \checkmark \\
MHGGNN & SMILES & 48 & 0.980 &  \\
PDV & Mol2Vec & 48 & 0.979 &  \\
ECFP4 & Rand.~SMILES & 16 & 0.979 & \checkmark \\
PDV & SMILES & 66 & 0.978 &  \\
MHGGNN & PDV (binary) & 54 & 0.972 & \checkmark \\
Mol2Vec & SNS & 44 & 0.969 & \checkmark \\
Mol2Vec & SMILES & 48 & 0.966 &  \\
ECFP4 & PDV (binary) & 78 & 0.965 & \checkmark \\
Morgan & PDV (binary) & 60 & 0.962 & \checkmark \\
PDV (binary) & SNS & 56 & 0.961 & \checkmark \\
Mol2Vec & PDV (binary) & 54 & 0.956 & \checkmark \\
MHGGNN & Rand.~SMILES & 16 & 0.956 & \checkmark \\
PDV (binary) & SMILES & 72 & 0.941 & \checkmark \\
PDV (binary) & Rand.~SMILES & 16 & 0.929 & \checkmark \\
Mol2Vec & Rand.~SMILES & 16 & 0.891 & \checkmark \\
\bottomrule
\end{longtable}}


\addfile{Intraclass correlation coefficients for model pairs}
\label{app:noise:sec:model_icc}
% TODO: review, needs manual review
ICC(1,1) of the noise degradation slope profiles for the model pairs with ICC $\geq 0.5$, on the QM9 HOMO--LUMO gap.
% Thesis appendix item: model_icc
% Seeded on 2026-09-27 from qsar_qm_models/supplementary/items/model_icc.tex by thesis_appendix/seed_from_paper.py.
% From here on this copy is edited independently of the paper. Rules: thesis_appendix/README.md.
% The lines below, down to the first blank line, are the PAPER copy's header, kept for
% provenance. They describe the paper file, not this one: in particular a GENERATED
% notice there does not apply here, and this copy may be edited by hand.
% Supplementary item: model_icc
% Split out of additional_files.tex on 2026-09-27, where it was Additional file 4 (ICC(1,1) (was Supp Table S5)).
% Content unchanged apart from: the heading line removed (the build writes it), the old
% 'Table S*/Figure S*' caption prefix replaced by \SuppFile{}, and the label renamed to supp:model_icc.
% Rules: supplementary/README.md.

\begin{table}[H]
\centering
\caption{\textbf{Intraclass correlation coefficients (ICC(1,1)) for model pairs with ICC $\geq 0.5$, assessing within-family consistency of noise degradation slope (NDS) profiles on the QM9 HOMO--LUMO gap.}
ICC(1,1) measures the proportion of total variance attributable to between-configuration differences; values near 1.0 indicate that two models rank configurations almost identically. Lower (more negative) NDS indicates faster degradation under noise. The Family column groups each pair: ``Tree'' (both tree-based), ``$\alpha$/$\beta$'' (architecture-pair partners such as BNN-$\alpha$/BNN-$\beta$), ``$\alpha$'' (both share the $\alpha$ neural-network backbone), or ``Cross'' (cross-family pair).}
\label{app:noise:model_icc}
\small
\begin{tabular}{llcccc}
\toprule
\textbf{Model A} & \textbf{Model B} & \textbf{$n$} & \textbf{ICC(1,1)} & \textbf{Mean $|\Delta$NDS$|$} & \textbf{Family} \\
\midrule
LightGBM & XGBoost & 8 & 0.944 & 0.0060 & Tree \\
BNN-$\alpha$ & BNN-$\beta$ & 6 & 0.846 & 0.0101 & $\alpha$/$\beta$ \\
QRF & RF & 7 & 0.810 & 0.0216 & Tree \\
VBLL-$\alpha$ & VBLL-$\beta$ & 4 & 0.782 & 0.0108 & $\alpha$/$\beta$ \\
QRF & XGBoost & 7 & 0.754 & 0.0157 & Tree \\
LightGBM & QRF & 7 & 0.743 & 0.0133 & Tree \\
VBLL-$\alpha$ & LightGBM & 4 & 0.698 & 0.0115 & Cross \\
VBLL-$\beta$ & SVM & 4 & 0.689 & 0.0133 & Cross \\
RF & XGBoost & 8 & 0.678 & 0.0212 & Tree \\
VBLL-$\alpha$ & XGBoost & 4 & 0.656 & 0.0133 & Cross \\
BNN-$\alpha$ & VBLL-$\alpha$ & 4 & 0.604 & 0.0134 & $\alpha$ \\
LightGBM & RF & 8 & 0.586 & 0.0244 & Tree \\
LightGBM & VBLL-$\beta$ & 4 & 0.577 & 0.0174 & Cross \\
LightGBM & SVM & 8 & 0.561 & 0.0117 & Cross \\
BNN-$\alpha$ & VBLL-$\beta$ & 4 & 0.552 & 0.0123 & $\alpha$/$\beta$ \\
VBLL-$\beta$ & XGBoost & 4 & 0.544 & 0.0195 & Cross \\
VBLL-$\alpha$ & QRF & 4 & 0.540 & 0.0229 & Cross \\
VBLL-$\alpha$ & SVM & 4 & 0.522 & 0.0124 & Cross \\
Gauche & VBLL-$\beta$ & 3 & 0.519 & 0.0113 & Cross \\
\bottomrule
\end{tabular}
\end{table}


\addfile{Global overview of noise robustness on ECFP4}
\label{app:noise:sec:ecfp4_overview}
% TODO: review, needs manual review
Noise robustness of every model on the ECFP4 representation, on the QM9 HOMO--LUMO gap.
% Thesis appendix item: ecfp4_overview
% Seeded on 2026-09-27 from qsar_qm_models/supplementary/items/ecfp4_overview.tex by thesis_appendix/seed_from_paper.py.
% From here on this copy is edited independently of the paper. Rules: thesis_appendix/README.md.
% The lines below, down to the first blank line, are the PAPER copy's header, kept for
% provenance. They describe the paper file, not this one: in particular a GENERATED
% notice there does not apply here, and this copy may be edited by hand.
% Supplementary item: ecfp4_overview
% Split out of additional_files.tex on 2026-09-27, where it was Additional file 6 (ECFP4 global overview (was Supp Figure S1)).
% Content unchanged apart from: the heading line removed (the build writes it), the old
% 'Table S*/Figure S*' caption prefix replaced by \SuppFile{}, and the label renamed to supp:ecfp4_overview.
% Rules: supplementary/README.md.

\begin{figure}[H]
\centering
\includegraphics[width=\textwidth,height=0.78\textheight,keepaspectratio]{fig1_supp_ecfp4_overview.png}
\caption{\textbf{Global overview of noise robustness on the ECFP4 representation, evaluated on QM9 HOMO--LUMO gap.}
a) Performance degradation curves (R$^2$ versus noise level $\sigma$) under Gaussian noise. b) Noise degradation slope (NDS) heatmap across model architectures and noise strategies; lower (more negative) NDS indicates faster degradation under noise.}
\label{app:noise:ecfp4_overview}
\end{figure}


\addfile{Clean $R^2$ against robustness on PDV under Gaussian noise}
\label{app:noise:sec:clean_r2_vs_robustness}
% TODO: review, needs manual review
Clean-label $R^2$ against the noise degradation slope for every model on PDV under Gaussian noise, on the QM9 HOMO--LUMO gap.
% Thesis appendix item: clean_r2_vs_robustness
% Seeded on 2026-09-27 from qsar_qm_models/supplementary/items/clean_r2_vs_robustness.tex by thesis_appendix/seed_from_paper.py.
% From here on this copy is edited independently of the paper. Rules: thesis_appendix/README.md.
% The lines below, down to the first blank line, are the PAPER copy's header, kept for
% provenance. They describe the paper file, not this one: in particular a GENERATED
% notice there does not apply here, and this copy may be edited by hand.
% Supplementary item: clean_r2_vs_robustness
% Split out of additional_files.tex on 2026-09-27, where it was Additional file 7 (Baseline R^2 vs NDS scatter (was Supp Figure S2)).
% Content unchanged apart from: the heading line removed (the build writes it), the old
% 'Table S*/Figure S*' caption prefix replaced by \SuppFile{}, and the label renamed to supp:clean_r2_vs_robustness.
% Rules: supplementary/README.md.

\begin{figure}[H]
\centering
\includegraphics[width=\textwidth]{fig3_ranking_consistency.png}
\caption{\textbf{Baseline R$^2$ versus noise degradation slope (NDS) scatter for all models on the physicochemical descriptor vector (PDV) representation under Gaussian noise, on the QM9 HOMO--LUMO gap.}
Lower (more negative) NDS indicates faster performance loss under noise.}
\label{app:noise:clean_r2_vs_robustness}
\end{figure}


\addfile{Bayesian transformation effects by representation}
\label{app:noise:sec:bayesian_transformation_tests}
% TODO: review, needs manual review
The effect of each Bayesian transformation of the neural networks on the noise degradation slope, by representation, on the QM9 HOMO--LUMO gap.
% Thesis appendix item: bayesian_transformation_tests
% Seeded on 2026-09-27 from qsar_qm_models/supplementary/items/bayesian_transformation_tests.tex by thesis_appendix/seed_from_paper.py.
% From here on this copy is edited independently of the paper. Rules: thesis_appendix/README.md.
% The lines below, down to the first blank line, are the PAPER copy's header, kept for
% provenance. They describe the paper file, not this one: in particular a GENERATED
% notice there does not apply here, and this copy may be edited by hand.
% Supplementary item: bayesian_transformation_tests
% Split out of additional_files.tex on 2026-09-27, where it was Additional file 8 (NN Bayesian transformation effects (was Supp Table S6)).
% Content unchanged apart from: the heading line removed (the build writes it), the old
% 'Table S*/Figure S*' caption prefix replaced by \SuppFile{}, and the label renamed to supp:bayesian_transformation_tests.
% Rules: supplementary/README.md.

\begin{table}[H]
\centering
\caption{\textbf{Neural network Bayesian transformation effects on noise degradation slope (NDS) by representation, on the QM9 HOMO--LUMO gap.}
Positive $\Delta$ NDS indicates the Bayesian variant is more robust (less negative slope) than its deterministic counterpart.}
\label{app:noise:bayesian_transformation_tests}
\small
\begin{tabular}{llrrrr}
\toprule
\textbf{Model} & \textbf{Representation} & \textbf{Baseline R$^2$} & \textbf{NN-$\beta$ NDS} & \textbf{BNN-$\beta$ NDS} & \textbf{$\Delta$ NDS} \\
\midrule
NN-$\beta$ & ECFP4 & 0.791 & $-0.464$ & $-0.388$ & $+0.076$ \\
NN-$\beta$ & PDV & 0.858 & $-0.487$ & $-0.382$ & $+0.106$ \\
NN-$\beta$ & SMILES & 0.721 & $-0.558$ & $-0.365$ & $+0.193$ \\
\bottomrule
\end{tabular}
\end{table}


\addfile{Uncertainty metrics for probabilistic models on ECFP4}
\label{app:noise:sec:uncertainty_metrics_ecfp4}
% TODO: review, needs manual review
Uncertainty metrics of the probabilistic models on ECFP4 under Gaussian noise, on the QM9 HOMO--LUMO gap.
% Thesis appendix item: uncertainty_metrics_ecfp4
% Seeded on 2026-09-27 from qsar_qm_models/supplementary/items/uncertainty_metrics_ecfp4.tex by thesis_appendix/seed_from_paper.py.
% From here on this copy is edited independently of the paper. Rules: thesis_appendix/README.md.
% The lines below, down to the first blank line, are the PAPER copy's header, kept for
% provenance. They describe the paper file, not this one: in particular a GENERATED
% notice there does not apply here, and this copy may be edited by hand.
% Supplementary item: uncertainty_metrics_ecfp4
% Split out of additional_files.tex on 2026-09-27, where it was Additional file 9 (Uncertainty metrics ECFP4 (was Supp Table S8)).
% Content unchanged apart from: the heading line removed (the build writes it), the old
% 'Table S*/Figure S*' caption prefix replaced by \SuppFile{}, and the label renamed to supp:uncertainty_metrics_ecfp4.
% Rules: supplementary/README.md.

\begin{table}[H]
\centering
\caption{\textbf{Uncertainty quantification metrics for probabilistic models on the ECFP4 representation (Gaussian noise strategy) on the QM9 HOMO--LUMO gap.}
Unc-Error $\rho$: Spearman correlation between predicted uncertainty and absolute prediction error. Unc-Noise $\rho$: Spearman correlation between predicted uncertainty and injected noise level $\sigma$. ECE: expected calibration error (lower = better-calibrated). Coverage targets: 68\% ($1\sigma$), 95\% ($2\sigma$).}
\label{app:noise:uncertainty_metrics_ecfp4}
\small
\begin{tabular}{lrrrrr}
\toprule
\textbf{Model} & \textbf{Unc-Error $\rho$} & \textbf{Unc-Noise $\rho$} & \textbf{ECE} & \textbf{Cov.\ $1\sigma$} & \textbf{Cov.\ $2\sigma$} \\
\midrule
\textbf{QRF}     & \textbf{0.25} & 0.27 & 0.19 & 74.7\% & 94.6\% \\
\textbf{NGBoost} & 0.22 & \textbf{0.37} & \textbf{0.11} & 69.2\% & 94.2\% \\
GP               & 0.18 & 0.34 & 0.23 & 78.6\% & 96.0\% \\
BNN-$\beta$      & 0.20 & 0.37 & 0.29 & 78.6\% & 95.2\% \\
BNN-$\alpha$     & 0.15 & 0.32 & 0.24 & 76.6\% & 95.1\% \\
VBLL-$\beta$     & 0.19 & 0.34 & 0.19 & 73.9\% & 95.5\% \\
VBLL-$\alpha$    & 0.11 & 0.31 & 0.22 & 74.5\% & 95.1\% \\
\bottomrule
\end{tabular}
\end{table}


\addfile{Variance decomposition on the three assay datasets}
\label{app:noise:sec:assay_variance_decomposition}
% TODO: review, needs manual review
ANOVA variance decomposition ($\eta^2$, in per cent) on LogD, Caco-2 efflux and hERG $K_i$.
% Thesis appendix item: assay_variance_decomposition
% Seeded on 2026-09-27 from qsar_qm_models/supplementary/items/assay_variance_decomposition.tex by thesis_appendix/seed_from_paper.py.
% From here on this copy is edited independently of the paper. Rules: thesis_appendix/README.md.
% The lines below, down to the first blank line, are the PAPER copy's header, kept for
% provenance. They describe the paper file, not this one: in particular a GENERATED
% notice there does not apply here, and this copy may be edited by hand.
% Supplementary item: assay_variance_decomposition
% Split out of additional_files.tex on 2026-09-27, where it was Additional file 10 (Validation ANOVA (was Supp Figure S3)).
% Content unchanged apart from: the heading line removed (the build writes it), the old
% 'Table S*/Figure S*' caption prefix replaced by \SuppFile{}, and the label renamed to supp:assay_variance_decomposition.
% Rules: supplementary/README.md.

\begin{figure}[H]
\centering
\includegraphics[width=\textwidth]{fig_validation_anova.png}
\caption{\textbf{ANOVA variance decomposition ($\eta^2$, \%) on external validation datasets: a) LogD, b) Caco-2 Efflux, and c) hERG K$_i$.}
Model, representation, and interaction contributions to noise degradation slope (NDS) variance for each dataset; lower (more negative) NDS indicates faster performance loss under noise.}
\label{app:noise:assay_variance_decomposition}
\end{figure}


\addfile{Random forest against quantile regression forest on the three assay datasets}
\label{app:noise:sec:rf_vs_qrf_assay}
% TODO: review, needs manual review
The noise degradation slope of random forest and of quantile regression forest on LogD, Caco-2 efflux and hERG $K_i$.
% Thesis appendix item: rf_vs_qrf_assay
% Seeded on 2026-09-27 from qsar_qm_models/supplementary/items/rf_vs_qrf_assay.tex by thesis_appendix/seed_from_paper.py.
% From here on this copy is edited independently of the paper. Rules: thesis_appendix/README.md.
% The lines below, down to the first blank line, are the PAPER copy's header, kept for
% provenance. They describe the paper file, not this one: in particular a GENERATED
% notice there does not apply here, and this copy may be edited by hand.
% Supplementary item: rf_vs_qrf_assay
% Split out of additional_files.tex on 2026-09-27, where it was Additional file 11 (RF vs QRF on validation datasets (was Supp Table S7)).
% Content unchanged apart from: the heading line removed (the build writes it), the old
% 'Table S*/Figure S*' caption prefix replaced by \SuppFile{}, and the label renamed to supp:rf_vs_qrf_assay.
% Rules: supplementary/README.md.

\begin{table}[H]
\centering
\caption{\textbf{Random forest (RF) vs quantile regression forest (QRF) noise degradation slope (NDS) on the three external validation datasets (LogD, Caco-2 Efflux, hERG K$_i$).}
Lower (more negative) NDS indicates faster degradation under noise.}
\label{app:noise:rf_vs_qrf_assay}
\small
\begin{tabular}{lrrr}
\toprule
\textbf{Dataset} & \textbf{RF NDS} & \textbf{QRF NDS} & \textbf{$\Delta$ NDS} \\
\midrule
ChEMBL-hERG-Ki     & $-0.231$ & $-0.359$ & $-0.128$ \\
OpenADMET-Caco2    & $-0.534$ & $-0.627$ & $-0.094$ \\
OpenADMET-LogD     & $-0.101$ & $-0.185$ & $-0.083$ \\
\bottomrule
\end{tabular}
\end{table}


\section{The flaws of the similarity assumption}\label{app:add:flaws}
% ---- NOTES: The flaws of the similarity assumption ----
% Takeaway: the per-dataset, per-representation counts behind Chapter 4's claims, in four tables.
% - Table, constant features. One row is one dataset and one representation. Columns: total features; constant features; near-constant features.
% - Table, molecules a representation cannot tell apart. One row is one representation and one dataset. Columns: molecules in groups with an identical feature vector; how many of those differ only in stereochemistry; lowest RMSE reachable when every group gets one prediction; largest label spread within a group.
% - Table, distance against property difference. One row is one dataset and one representation. Columns: rank correlation between pairwise distance and absolute property difference; score of predicting each molecule's activity from its five nearest neighbours.
% - Table, activity cliffs. One row is one dataset and one representation. Column: molecules in at least one activity cliff.
% Result needed: all four tables, from the full runs.
% ---- END NOTES ----
\todo[inline, color=orange!40]{Additional files for Chapter~\ref{chap:flaws}: all four need the full runs.}

\section{Building adaptive hybrid representations}\label{app:add:method}\label{app:searchspace}
% ---- NOTES: Building adaptive hybrid representations ----
% Takeaway: these files are what let a reader rebuild any hybrid in the thesis. Chapter 5 gives the reasons; this section holds only the reference tables.
% - Change: the design-experiment file's allocation row gives only counts. Add the reason the global pool with squared weights was kept over the other pooled candidates, once Chapter 5 states it.
% Result needed: the best representation count per (downstream model, budget) from the design run, for the design-experiment file's representation-count row.
% Result needed: per dataset and downstream model, per-fold score of the best single representation against the hybrid; the winning configuration; how often each setting value won.
% Result needed: augmentation per dataset: how often no extra features were chosen; per-fold base against base plus extras, uncapped and capped; which extra feature sets were kept; decoy features admitted.
% Result needed: per dataset and budget, held-out score against the score the inner search reported, and how stable the selected features are across folds.
% Result needed: hybrid minus each comparator (concatenation, one model per representation, PCA, mRMR, recursive feature elimination, random features) at each budget.
% Result needed: the smaller comparator, control and ablation experiments, one table each, if they leave Chapter 5.
% Result needed: per task, outcome of the transfer, antibody-scan and separate-site experiments, with which representations the selected features came from.
% Result needed: hours and model fits per dataset at each tier, with the hardware.
% ---- END NOTES ----

\addfile{Settings of every model, as downstream model and as internal evaluator}
\label{app:add:model_settings}
% TODO: review, needs manual review
One row is one model. The middle two columns give its settings when it scores a hybrid (downstream model) and when it scores one representation inside the search (internal evaluator). The models fitted only to compute feature importance are described in Chapter~\ref{chap:method}.

% TODO: review, needs manual review (whole table and caption, written from the current code)
{\small
\begin{longtable}{@{}p{1.9cm}p{4.6cm}p{3.6cm}p{1.6cm}p{2.3cm}@{}}
\caption{Settings of every model. $p$ is the number of features. ``Molecules'' is the number of training molecules the internal evaluator is fitted on; above it, a random subset of that size is drawn with a fixed seed. Settings not listed are the library's defaults. Every setting is the same on every dataset. The Gaussian process runs only on the datasets named in Chapter~\ref{chap:method}, ``Models and metrics''.}
\label{tab:add:model_settings}\\
\toprule
Model & As downstream model & As internal evaluator & Molecules & Library; source \\
\midrule
\endfirsthead
\toprule
Model & As downstream model & As internal evaluator & Molecules & Library; source \\
\midrule
\endhead
Random forest & 100 trees, no depth limit; $\sqrt{p}$ features tried at each split & 50 trees, depth at most 10; all $p$ features tried at each split on regression, $\sqrt{p}$ on classification & all & scikit-learn \citep{scikit-learn}; \citet{Breiman2001} \\
\addlinespace
LightGBM & 100 trees & 50 trees, depth at most 4 & all & lightgbm; \citet{ke2017lightgbm} \\
\addlinespace
XGBoost & 100 trees; half of the features sampled for each tree & 50 trees, depth at most 4 & all & xgboost; \citet{Chen2016} \\
\addlinespace
SVM & RBF kernel, $C = 1$. Not fitted above 10{,}000 molecules, or above 5{,}000 under uncapped augmentation & RBF kernel, $C = 1$ & at most 500 & scikit-learn; \citet{Vapnik1995} \\
\addlinespace
Gaussian process & Regression: RBF plus white-noise kernel, labels normalised, 3 optimiser restarts. Classification: RBF kernel. Not fitted above 1{,}000 molecules & Regression: scaled RBF kernel, fixed noise variance 0.01, no restarts. Classification: scaled RBF kernel, no restarts & at most 500 & scikit-learn; \citet{Rasmussen2006} \\
\addlinespace
DNN & Hidden layers of 256, 128 and 64 units, ReLU, dropout 0.2 \citep{Srivastava2014}; Adam \citep{Kingma2015} at learning rate 0.001, halved after 5 epochs without improvement; batch size 256; up to 200 epochs; stops after 20 epochs without improvement on 20\% of the training molecules that share no group with the rest, using the outer split's grouping; labels z-normalised & The same, except that the early-stopping 20\% is drawn at random & all & PyTorch \citep{Paszke2019} \\
\bottomrule
\end{longtable}
}

\addfile{Downstream models left out, by dataset and augmentation setting}
\label{app:add:models_left_out}
% TODO: review, needs manual review
One row is one rule that leaves a downstream model out. Only the 16 datasets of Table~\ref{tab:datasets} are listed; a dataset and setting not named in any row runs every downstream model in Table~\ref{tab:search_tiers}.

% TODO: review, needs manual review (whole table and caption, written from the current code)
{\small
\begin{longtable}{@{}p{2.4cm}p{3.0cm}p{8.6cm}@{}}
\caption{Downstream models left out, by dataset and augmentation setting. Every rule was set so that a search finishes within the time limit of one job; none was set from scores. hERG $K_i$ drops no model under any setting.}
\label{tab:add:models_left_out}\\
\toprule
Model left out & Augmentation setting & Datasets \\
\midrule
\endfirsthead
\toprule
Model left out & Augmentation setting & Datasets \\
\midrule
\endhead
XGBoost & all three & H3K4me3, TATA promoter \\
SVM & all three & datasets above 10{,}000 molecules: H3K4me3 \\
SVM & uncapped & datasets above 5{,}000 molecules: QM9, LogD$^{a}$, TATA promoter$^{a}$ \\
SVM & uncapped & Koenig expression \\
SVM & capped & Caco-2 efflux, LogD, rat liver microsome clearance, MDR1--MDCK efflux, Buchwald--Hartwig plates \\
Random forest and SVM & capped & QM9 \\
Random forest & uncapped & QM9, LogD \\
Gaussian process & all three & every dataset except humanness (Chapter~\ref{chap:method}, ``Models and metrics'') \\
\bottomrule
\end{longtable}
{\footnotesize $^{a}$Counted on the whole dataset: LogD has 5{,}028 molecules and TATA promoter 5{,}274. A search run inside one outer fold is given about 80\% of them and keeps the SVM.}
}

\addfile{Size of the search}
\label{app:add:search_size}
% TODO: review, needs manual review
One row is one downstream model. The counts are per dataset, objective and augmentation setting, with adaptive pruning off.

% TODO: review, needs manual review (whole table and caption)
\begin{table}[htbp]
\centering
\small
\caption{Number of configurations searched per dataset, objective and augmentation setting. A configuration is one choice of downstream model, internal evaluator, allocation method, importance method, representation count and budget. At the full tier each (downstream model, internal evaluator) pair runs 5 allocation methods $\times$ 3 importance methods $\times$ 4 representation counts $\times$ 8 budgets $= 480$ configurations; at the fast tier it runs $3 \times 1 \times 2 \times 8 = 48$ (Table~\ref{tab:search_tiers}). ``Matched'' is counted once where it repeats an evaluator already in the list. Each configuration is scored on 5 folds, so the full tier fits a downstream model 40{,}800 times (50{,}400 on humanness) and the fast tier 4{,}080 times. Where a dataset leaves a downstream model out (Additional file~\ref{app:add:models_left_out}), its pairs drop out: a fast search on H3K4me3 has 9 pairs, so 432 configurations.}
\label{tab:add:search_size}
\begin{tabular}{@{}lp{4.2cm}rr@{}}
\toprule
Downstream model & Internal evaluators & Full tier & Fast tier \\
\midrule
Random forest & 3 (matched is random forest) & 1{,}440 & 144 \\
LightGBM & 3 (matched is LightGBM) & 1{,}440 & 144 \\
DNN & 3 (matched is DNN) & 1{,}440 & 144 \\
XGBoost & 4 & 1{,}920 & 192 \\
SVM & 4 & 1{,}920 & 192 \\
Gaussian process (humanness only) & 4 & 1{,}920 & not run \\
\midrule
Total & 17 pairs (21 on humanness) & 8{,}160 (10{,}080 on humanness) & 816 \\
\bottomrule
\end{tabular}
\end{table}

\addfile{Design experiments behind the search settings}
\label{app:add:design_experiments}
% TODO: review, needs manual review
One row is one setting of the search, or one rule that was tested and then fixed instead of searched. Every number comes from preliminary design runs made before the main experiments; they are not results of the thesis, and several used other budgets than the final search.

% TODO: review, needs manual review (whole table and caption; numbers from preliminary design runs)
{\footnotesize
\begin{longtable}{@{}p{1.9cm}p{2.5cm}p{2.7cm}p{4.4cm}p{2.5cm}@{}}
\caption{Design experiments behind the search settings. Preliminary design runs, not the thesis's full runs.}
\label{tab:add:design_experiments}\\
\toprule
Setting & Values compared & Datasets, folds, budgets & What was measured, and the result & Values kept (full / fast) \\
\midrule
\endfirsthead
\toprule
Setting & Values compared & Datasets, folds, budgets & What was measured, and the result & Values kept (full / fast) \\
\midrule
\endhead
Downstream model & Nine models, among them random forest, LightGBM, XGBoost, CatBoost, DNN, $k$-nearest neighbours and a Gaussian process, each on each of 10 base representations at full width & Buchwald--Hartwig yield, Caco-2 efflux, Koenig expression, LogD, breast cancer PAM50 subtype; 5 folds & Best held-out score over the 10 representations, and seconds per fit. The DNN scored best on 4 of the 5 datasets and third on Caco-2 (0.2936 against 0.2851 RMSE, log$_{10}$ efflux ratio) & Full: random forest, LightGBM, XGBoost, SVM, DNN, and the Gaussian process on humanness / fast: the same five \\
\addlinespace
Budget & None compared & -- & -- & 10, 25, 50, 100, 250, 500, 1{,}000 and 2{,}000 features at both tiers \\
\addlinespace
Representation count & 2, 3, 5 and all 10 & 9 datasets, 3 tree-ensemble downstream models, budgets 500 and 1{,}000 features; 5 folds & The best count for each (downstream model, budget) combination & Full: 2, 3, 5, all / fast: 3, all \\
\addlinespace
Allocation method & 8 candidates: the five searched and 3 further pooled candidates, each at every representation count & The same 9 datasets, models and budgets: 54 (dataset, model, budget) combinations & Each candidate was scored at its best representation count, and the number of combinations in which it was in the top 3 of the 8 was counted. Fixed shares: linear 30, squared 27, rank decay 26, softmax 25; global pool with squared weights 14 & Full: the five described in Chapter~\ref{chap:method} / fast: linear, squared, rank decay \\
\addlinespace
Importance method, agreement & 8 methods & LogD, hERG $K_i$, Buchwald--Hartwig yield, nanobody thermostability, antibody FcRn retention, breast cancer PAM50 subtype, lung adenocarcinoma vs squamous; each method fitted once on the whole dataset & Spearman correlation between two methods' scores over every feature. Kept pairs: LightGBM gain and TreeSHAP 0.20--0.64, LightGBM gain and DeepLIFT 0.21--0.67, TreeSHAP and DeepLIFT 0.56--0.99 across the 7 datasets & One method kept from each group whose rankings agree above 0.95 \\
\addlinespace
Importance method, score & 28 methods, among them three controls: feature variance, Laplacian score and random ranking & Caco-2 efflux; 3 folds; budget 100 features; random forest and LightGBM downstream & RMSE of the best single representation minus RMSE of the hybrid (log$_{10}$ efflux ratio; positive means the hybrid has the lower error): DeepLIFT $+0.0144$, LightGBM gain $+0.0100$, TreeSHAP $+0.0067$, random ranking $-0.0118$ & Full: LightGBM gain, TreeSHAP, DeepLIFT / fast: LightGBM gain \\
\addlinespace
Internal evaluator settings & 20 hyperparameter settings from 5 model families & Buchwald--Hartwig yield, Caco-2 efflux, Koenig expression; 3 folds; budget 100 features & Spread of mean scores across all 20 settings: 0.0253 $R^2$, 0.0179 $R^2$ and 0.0085 Spearman. Median standard deviation across the 3 folds for one setting, allocation method and downstream model: 0.0462 $R^2$, 0.0693 $R^2$ and 0.0186 Spearman & Matched, random forest, LightGBM and DNN, at fixed settings, both tiers \\
\addlinespace
Rescaling of each representation's score (a rule) & 5 rescalings & 9 datasets; 5 folds; budgets 50 and 300 features & Whether any rescaling beat the one in use by more than that rule's own spread across folds: in 0 of 90 (dataset, allocation method, budget) combinations & One fixed rule per metric (Chapter~\ref{chap:method}) \\
\bottomrule
\end{longtable}
{\footnotesize The representation-count and allocation rows come from one design run on 10 datasets, 3 downstream models and 2 budgets (60 combinations); QM9 is left out, which gives the 54.}
}

\addfile{Fields a blueprint saves}
\label{app:add:blueprint_fields}
% TODO: review, needs manual review
One row is one field. The first block is the recipe that rebuilds a hybrid; the second is kept for inspection and is not read by a rebuild.

% TODO: review, needs manual review (whole table and caption, written from the current code)
{\small
\begin{longtable}{@{}p{4.6cm}p{9.4cm}@{}}
\caption{Fields a blueprint saves. The stored score of each representation is its best score on its own across the downstream models, over every molecule the search ran on. It is not the internal-evaluator score that splits the budget, which a rebuild computes again on the molecules it is given. Three further fields are read only by allocation or augmentation routes outside the thesis's search and are fixed: a minimum improvement a whole block of extra features must show (none), a number of features taken from each source (20), and a step size (5).}
\label{tab:add:blueprint_fields}\\
\toprule
Field & What it holds \\
\midrule
\endfirsthead
\toprule
Field & What it holds \\
\midrule
\endhead
\multicolumn{2}{@{}l}{\textit{The recipe}} \\
\midrule
Representations & The representations offered to the search \\
Representation count & How many of them the configuration used; empty when it used all \\
Allocation method & One of the five described in Chapter~\ref{chap:method} \\
Importance method & LightGBM gain, TreeSHAP or DeepLIFT \\
Budget & Number of features in the hybrid \\
Internal evaluator & Matched, random forest, LightGBM or DNN \\
Objective & Predictive, or predicting each molecule's activity from its five nearest neighbours \\
Metric & The dataset's metric, and whether lower is better \\
Task type & Regression or classification \\
Folds and seed & The number of folds, and the random seed \\
\midrule
\multicolumn{2}{@{}l}{\textit{The recipe: augmentation}} \\
\midrule
Augmentation setting & None, uncapped or capped \\
Candidate ranking & How candidate extra features are ranked \\
Candidate list length & 50 features \\
Cap per source & Optional; the most candidates taken from one source \\
Margin & How much better than adding none a block of extra features must score \\
Fraction of the budget & Optional; used instead of the fixed block sizes \\
\midrule
\multicolumn{2}{@{}l}{\textit{Kept for inspection, not read by a rebuild}} \\
\midrule
Representation scores & Each representation's best score on its own across the downstream models, sign flipped for metrics where lower is better \\
Features per representation & The number of features each representation got in each fold \\
Extra features kept & The extra features augmentation kept in each fold, by source \\
\bottomrule
\end{longtable}
}

\todo[inline, color=orange!40]{Additional files for Chapter~\ref{chap:method}: the per-dataset results, augmentation, held-out against inner-search scores, comparators, smaller experiments, transfer experiments and compute need the full runs.}

\section{Hybrid representations for similarity and label noise}\label{app:add:similarity}
% ---- NOTES: Hybrid representations for similarity and label noise ----
% Takeaway: the full tables behind Chapter 6's summaries, and the settings of the twelve LNL methods.
% - Table, finding injected noisy labels. One row is one noise score. Columns: average precision under outlier noise at 1\%, 5\% and 10\%, grouped-wider noise and censoring, in each feature space; chance level in the header row.
% - The feature spaces: each model's top two single representations, the same features weighted by importance, forest proximity, the hybrid, and the hybrid ranked on the neighbour score.
% - Table, finding real suspect labels. One row is one noise score and feature space. Columns: average precision against ChEMBL curator flags on hERG IC50, per record and per molecule, with chance level; Spearman correlation with the spread between laboratories on the BindingDB molecules with repeat measurements.
% - Table, estimated noise size. One row is one dataset and feature space. Columns: the estimate from repeat measurements (with the measured variance between repeats where the dataset keeps them); the estimate from near-duplicate pairs; the gamma-test estimate; the model's held-out mean squared error. All in label units squared.
% Result needed: the three tables above, from the full runs.
% Result needed: per dataset and fold, importance-weighted against unweighted features for predicting each molecule's activity from its five nearest neighbours; the hybrid searched for that objective against the ordinary hybrid and the two best single representations.
% Result needed: censoring: share of the lost score recovered, per treatment, model, feature space and clipping level; deleting clipped molecules against keeping them.
% Result needed: the survey of twelve LNL methods: per noise level, the methods that beat the plain network on each input.
% Result needed: cleaning real labels: folds won per cleaning method, noise score and feature space.
% Result needed: per noise type, whether any LNL method, noise mitigation or noise-robust training, improved on the untreated model, on each feature space.
% ---- END NOTES ----

\addfile{Settings of the twelve LNL methods and the plain network they share}
\label{app:add:lnl_settings}
% TODO: review, needs manual review
One row is one method; the first row is the plain network every method starts from. The kind of each method follows Chapter~\ref{chap:similarity}.

% TODO: review, needs manual review (whole table and caption, written from the current code)
{\small
\begin{longtable}{@{}p{2.9cm}p{2.0cm}p{7.3cm}p{1.6cm}@{}}
\caption{Settings of the twelve LNL methods. On classification datasets the squared-error loss is replaced by cross-entropy.}
\label{tab:add:lnl_settings}\\
\toprule
Method & Kind & Settings & Labels \\
\midrule
\endfirsthead
\toprule
Method & Kind & Settings & Labels \\
\midrule
\endhead
Plain network & -- & Hidden layers of 256, 128 and 64 units, ReLU, dropout 0.2; Adam at learning rate 0.001, halved after 5 epochs without improvement; batch size 256; gradient norm clipped at 1; up to 200 epochs; stops after 20 epochs without improvement on 20\% of the training molecules held out by scaffold group (drawn at random after a method has removed molecules); labels z-normalised & both \\
\midrule
Small-loss selection & noise mitigation & 10 warm-up epochs on all molecules; keeps the 70\% with the smallest loss and fits a new plain network on them & both \\
Confident learning & noise mitigation & 3 plain networks; removes molecules whose error is above the 85th percentile and whose spread across the 3 networks is below the 25th percentile, then refits & both \\
Removal by distance to neighbours & noise mitigation & 10 nearest neighbours in the standardised features; score is the distance of a label from its neighbours' mean label, divided by their standard deviation (on classification, the share of neighbours with another label); removes molecules above the 85th percentile, then refits & both \\
Staged removal & noise mitigation & 3 stages; each removes the 10\% of the remaining molecules with the largest training loss & both \\
\midrule
Per-epoch down-weighting & noise-robust training & After each epoch, the weight of every molecule whose loss is above the median is multiplied by 0.95 & both \\
Curriculum by confidence & noise-robust training & 20 warm-up epochs, then trains on the 25\%, 50\%, 75\% and 100\% of molecules with the smallest loss, in stages & both \\
DivideMix & noise-robust training & Two networks, 10 warm-up epochs; a two-component Gaussian mixture on each network's squared errors splits the molecules & both \\
Meta-Weight-Net & noise-robust training & The weighting network is updated every 5 epochs at learning rate 0.001, on 20\% of the training molecules held out for it & both \\
MentorNet & noise-robust training & The mentor network reads each molecule's current loss, its last 5 losses and its mean distance to its 10 nearest neighbours & both \\
Contrast to Divide & noise-robust training & 5 perturbed copies of each molecule, Gaussian noise of standard deviation 0.01 on the features & both \\
\midrule
Heteroscedastic GP & noise-robust training & RBF kernel, 2 optimiser restarts, fitted on at most 1{,}000 training molecules; a small network predicts each molecule's noise variance & regression \\
Deep evidential regression & noise-robust training & Weight of the evidence penalty 0.01 & regression \\
\bottomrule
\end{longtable}
}

\todo[inline, color=orange!40]{Additional files for Chapter~\ref{chap:similarity}: finding injected and real noisy labels, estimated noise size and the other results need the full runs.}
